% Folder 142, batch 05 — archivists' pages 81–100.
% Pass: Opus 5.5 (claude-opus-5-5), 2026-10-02 — first pass, unchecked against the pages by a human.
%
% Fast working hand throughout; the formulas carry, much of the connecting
% prose does not. The author's own pagination (8 … 17) and his running
% equation numbers (34) … are recorded where they help follow references.
\documentclass[11pt,a4paper]{article}
\input{../preamble/grothendieck}

\folder{142}
\batch{5}
\pages{81}{100}
\foldertitle{Action de Galois sur Teichmüller : notes manuscrites (s.d.).}
\dating{[à partir de 1978]}
\watermark{Édition de démonstration}

\begin{document}

\page{81}
\note{la phrase qui ouvre la page vient de la page précédente, hors de ce lot ; la page porte, en haut, la suite des équations numérotées (34) et suivantes de l'auteur}

\ill{} de façons équivalentes, par les $g^{\alpha}_{x,y}$ ($\in \hat{\Pi}_{\alpha}$), où $\alpha \in A$, $x, y \in I_{\alpha}$ (cf. plus haut). Ceci rappelé,

\note{le passage qui suit, jusqu'à « commencé ! », est marqué d'un trait vertical dans la marge gauche}
on peut espérer de trouver des « formules » explicites, dépendant \struck{bien} sûr des \textit{propriétés arithmétiques précises des points particuliers $x$}, et \textit{caractéristiques} de $x$ (dans la situation Mordell--Weil--Néron), qui donnent $f_{\ell}(u)$ \textit{en termes} des paramètres précédents fixant $u$. C'est là un travail fascinant en perspective, que je n'ai pratiquement pas encore commencé !

À vrai dire, la dimension\uncertain{alité} via Mordell--Weil \ill{} d'un \struck{$x$} \add{$\in X_K$} \struck{\ill{}} correspond à une \ill{} dans $H^1(\Gamma, \hat{\pi}_1(x_0))$), \ill{} \struck{\ill{}} \add{\uncertain{connaît}} seulement \ill{} donc \ill{} dans \ill{} \uncertain{unicité} $=$ \ill{} l'élément correspondant dans
\[
  (34) \qquad (f_{x_0, x})_{\mathrm{ab}} \in H^1\bigl(\Gamma, \hat{\pi}_1(X, x_0)_{\mathrm{ab}}\bigr),
\]
et \ill{} de déterminer, par des « formules explicites », les $(f_{\ell})_{\mathrm{ab}}$ correspondant aux chemins $x_0 \to x$. On en vient à regarder la \struck{\ill{}} question des \ill{} par une VA \ill{} cf. d'icelle

\marginal{en biais dans la marge gauche, en face de (34) : C'est une question particulière, \ill{} plus \uncertain{intéressante} \ill{} une formule \ill{} \ill{} \ill{} \uncertain{en remplaçant} $X$ par \ill{}}

\page{82}
\note{p. 9 de l'auteur}
par un tore --- pratiquement, par une « jacobienne généralisée ») convenable de $X$. Mais il faut se rappeler que pour une telle variété (p.\,ex. une courbe elliptique, ou la variété $\mathbb{G}_m$ ($\mathbb{G}_m(\mathbb{C}) = \mathbb{C}^{*}$)), l'opération de $\Gamma$ sur le groupoïde $\Pi_1(X_{\bar K})$ \uncertain{n'est pas} fidèle, donc pour « repérer » un $u \in \Gamma$, il ne suffit pas de connaître des invariants du type $g_{\ell}(u)$ relatifs à une \textit{seule} \struck{même} $X$. Il faudra donc \ill{} introduire des « paramètres de repérage » pour $u$, en utilisant d'autres variétés algébriques que $X$ --- la plus naturelle étant
\[
  (35) \qquad U_{0,3} = \mathbb{P}^1 \setminus \{0, 1, \infty\}.
\]

Faisons maintenant varier $X$, en \struck{\ill{}} considérant un morphisme de $K$-schémas
\[
  (36) \qquad f : X \to Y ;
\]
il définit des morphismes de groupoïdes

\page{83}
\[
  (37) \qquad \Pi_1(X) \xrightarrow{\; f_{*} = \Pi_1(f) \;} \Pi_1(Y),
\]
\[
  (38) \qquad \hat{\Pi}_1(X) \xrightarrow{\; \hat{f} = \hat{\Pi}_1(f) \;} \hat{\Pi}_1(Y),
\]
dont le deuxième peut être déduit du premier par voie de \uncertain{complétion} profinie ; mais par contre sa définition \ill{} \ill{} en termes \add{des morphismes} \ill{} schémas sur le corps alg.\ clos $\mathbb{C}$. Par \ill{} schématique, \ill{} $X \to Y$ sur le \ill{} \ill{}

\ill{} (38) commute aux actions : l'action de $\Gamma$, \add{d'où}
\[
  (39) \qquad \hat{f}({}^{u}\ell) = {}^{u}(\hat{f}\ell)
\]
et si on se borne aux sous-groupoïdes induits sur $X_{\bar K}$, $Y_{\bar K}$, $(\hat{\Pi}_1 Y_{\bar K})$, \ill{} : l'action des \ill{} $\Gamma = \mathrm{Gal}\, \bar{K}/K$. Si on se \ill{} \ill{} $X_K$, d'où \ill{} de $f_{\ell}(u)$, $g_{\ell}(u)$ \ill{} pour un chemin $\ell$ dans $X \to$ \ill{} $Y$, il vient
\[
  (40) \qquad
  \begin{cases}
    \hat{f}\bigl(f_{\ell}(u)\bigr) = f_{\hat{f}(\ell)}(u) \\
    \hat{f}\bigl(g_{\ell}(u)\bigr) = g_{\hat{f}(\ell)}(u) .
  \end{cases}
\]

\marginal{en biais dans la marge gauche : Analogue formule \ill{} \ill{} $+$ \ill{} $T$ \ill{} $(T \to \ill{}$ \ill{} \ill{} \ill{}) : (40 bis) $g_{f(x), f(y)}(u) = \hat{f}\bigl(g_{x,y}(u)\bigr)$ \ill{}}

\page{84}
\note{p. 10 de l'auteur}
Un cas trivial extrême est celui où $f$ est un isomorphisme --- et \struck{\ill{}} celui où $X = Y$, donc où $f$ fait partie d'un groupe $\mathfrak{G}$ de $K$-isomorphismes de $X$ --- dont on peut considérer $\Pi_1(X)$ et $\Pi_1(X_{\bar K})$, $\Pi_1(X_K)$ comme des groupoïdes à opérateurs $\mathfrak{G}$, et $\hat{\Pi}_1(X)$ comme un groupoïde à opérateurs $\mathfrak{G} \times \Gamma$, $\hat{\Pi}_1(\bar{X}_{\bar K})$ comme un groupoïde \struck{\ill{}} \uncertain{analogue} \ill{} $\mathfrak{G} \times \Gamma$ opère. Cette fois-ci les groupes \ill{} \uncertain{trivialement} sur les \uncertain{sommets} de $\Pi_1(X_K)$, l'opération de $\Gamma$ est triviale, mais pas nécessairement celle de $\mathfrak{G}$.

Si on se donne un pt base, et \uncertain{plus} généralement un ensemble $I \subset X_K$ contenant $x_0$, l'opération de $\Gamma$ sur $\hat{\Pi}_1(X, I)$ est connue quand on connaît son opération sur $\hat{\pi}_1(X, x_0)$, et si de plus on connaît son opération sur les $\hat{\pi}_1(X, x_0)$-torseurs ! droits
\[
  (41) \qquad \boxed{\mathrm{Hom}_{\hat{\Pi}_1(X, I)}(x_0, x_i)}, \qquad i \in I,
\]
ou, ce qui revient au même,

\page{85}
\[
  (42) \qquad f_{\ell_i}(u), \qquad i \in I,\ u \in \Gamma,
\]
où pour tout $i \in I$, $\ell_i$ est un chemin de $\hat{\Pi}_1(X, I)$ de $x_0$ vers $x_i$, i.e.\ un élément de la partie « discrète » du torseur (41).

D'autre part, l'action de $\Gamma$ sur $\hat{\pi}_1(X, x_0)$ est connue quand on connaît son action sur des générateurs $\lambda_j$ ($j \in J$) de $\pi_1(X, x_0)$, i.e.\ si on connaît les
\[
  (43) \qquad f_{\lambda_j}(u) \ \text{ou}\ g_{\lambda_j}(u), \qquad j \in J,\ u \in \Gamma .
\]
\struck{\ill{}} Les données (42), (43) suffisent \ill{} : déterminent l'action de $\Gamma$ sur $\Pi_1(X, I)$ --- mais ces descriptions, \struck{et} \struck{dis\ill{}} ; il a \ill{} d'un \ill{} des \ill{} \ill{}, à la dissymétrie des choix arbitraires, le choix d'un \struck{origine} \add{point-base} $x_0$ \ill{} \uncertain{nommément} \ill{} \ill{} \ill{} --- ce qui semble artificiel, et à la \uncertain{longue} je crois impraticable, \struck{et} et tout particulièrement \ill{} des situations \ill{} un groupe de symétries $\mathfrak{G}$ \ill{} \ill{} \ill{} \struck{\ill{}} puis $I$ stable sous $\mathfrak{G}$ ---

\marginal{en biais dans la marge gauche, en face de (43) : \ill{} \ill{} principe, \ill{} \struck{\ill{}} \uncertain{dit} p.~8'}

\page{86}
\note{p. 11 de l'auteur}
alors que pratiquement jamais on ne peut trouver un $x \in X_K$ \struck{\ill{}} fixe sous $\mathfrak{G}$.

Avant de donner un principe d'approche pour des \ill{} « symétriques », tenant compte de la présence d'un groupe $\mathfrak{G}$ (cas d'un groupe \ill{} \ill{} \uncertain{pratiques}), je voudrais préciser la question de la relation entre $\Pi_1(X, I)$ ($I$ partie finie de $X_K$) et $\Pi_1(U_{0,3})$, \struck{\ill{}} dans l'optique des opérations de $\Gamma$. Pour \struck{fixer les} idées \uncertain{pour simplifier}, je prendrai le cas $X$ \uncertain{quasi-projective} lisse géom.\ connexe. Utilisant des th.\ du type Lefschetz, plus le résultat de \struck{Biel} Belyi, on \struck{\ill{}} doit pouvoir trouver un \ill{} de schémas sur $K$
\note{une insertion au-dessus de la ligne, peu lisible, semble ajouter « $K$ à $\Gamma$ finis » ; non retenue}

\begin{tikzcd}
  C \arrow[r, hook, "i"] \arrow[d, "f"'] & X \\
  U_{0,3} &
\end{tikzcd}

\note{diagramme (44) de l'auteur ; au-dessus de $C$ il écrit $I$ avec le signe $\cap$, soit $I \subset C$}

où $i$ induit un hom.\ \textit{surjectif} sur les groupes fondamentaux, et où $f$ est

\page{87}
un morphisme fini étale. Ainsi $\Pi_1(X, I)$ apparaît comme un groupoïde quotient de $\Pi_1(C, I)$, donc $\hat{\Pi}_1(X, I)$ quotient de $\hat{\Pi}_1(C, I)$, et où $\Pi_1(C, I)$ apparaît comme un \ill{} fini de $\Pi_1(U_{0,3}, I)$ --- plus \struck{précisément} \add{exactement}, \uncertain{prenant le saturé} $I' = f^{-1}f(I) = f^{-1}(I_0)$ de $I$, $\Pi_1(U_{0,3}; I)$ est un sous-groupoïde \add{plein} de $\Pi_1(U_{0,3}; I')$, qui lui-même est un \uncertain{revêtement} fini des groupoïdes de $\Pi_1(U_{0,3}; I_0)$ --- et \uncertain{idem} pour les complétés profinis.
\note{dans les deux dernières formules on attendrait $C$ plutôt que $U_{0,3}$ ; on transcrit ce que porte la page}

Ainsi, l'opération de $\Gamma$ sur $\hat{\Pi}_1(X)$ est connue quand on connaît l'opération sur $\hat{\Pi}_1(C, \tilde{I})$, qui est un \ill{} fini de $\hat{\Pi}_1(U_{0,3}; I_0)$. Si on connaît l'action de $\Gamma$ sur $\hat{\Pi}_1(U_{0,3}; I_0)$ lui-même, \struck{\ill{} doit \ill{} \ill{}}

\page{88}
\note{p. 12 de l'auteur}
$\hat{\Pi}_1(C, I')$ est déterminé \add{via (39) et (40)}, \ill{} l'\uncertain{homomorphisme} \ill{} : pour $x, y \in C$, l'application
\[
  \mathrm{Hom}_{\hat{\Pi}_1(C)}(x, y) \longrightarrow \mathrm{Hom}_{\hat{\Pi}_1(U_{0,3})}(x, y)
\]
est injective !

J'ai été un peu \uncertain{vite} dans ce qui précède ; car le saturé $I'$ \add{dans $C$} de $I$ n'est pas \uncertain{tous} \struck{\ill{}} contenu dans $C_K$ --- mais \uncertain{dans} $C_{K'}$ \ill{} : mais dès lors, rien ne \uncertain{oblige} à remplacer $I$ par $I'$ \struck{\ill{}} pour faire \ill{} : il reste vrai que la connaissance de l'action de $\Gamma$ sur $\hat{\Pi}_1(U_{0,3}, I_0)$ implique la connaissance de l'action de $\Gamma$ sur $\hat{\Pi}_1(X, I)$. Mais il n'est \ill{} que l'étude des relations entre $\hat{\Pi}_1(C; I)$ et $\hat{\Pi}_1(U_{0,3}, I_0)$, en tant que groupoïdes à groupe d'opérateurs $\Gamma$, est formellement équivalente en remplaçant $I$ par son saturé $I'$, si celui-ci n'est pas $\subset C_K$. (Au besoin, si on est \struck{\ill{}} \add{pas} \ill{} par des arguments qui \ill{} \uncertain{ballottent} \ill{} l'action de $\Gamma$, il y a

\page{89}
qui remplacerait $\Gamma$ par un sous-groupe ouvert convenable $\Gamma'$ \ldots)

En même temps, on cherchera des conditions nécessaires et suffisantes pour qu'un automorphisme $\Theta$ des groupoïdes $\Pi_1(U_{0,3}; I_0)$, laissant fixes les sommets i.e.\ les pts de $I_0$, provienne d'un $u \in \Gamma$ : il faut que pour tout revêtement $C$ de $U_{0,3}$, défini sur $K$, \struck{\ill{}} \add{posant} $I = f^{-1}(I_0) \subset C_K$ et \struck{supposant} \add{$I \to I_0$ surj.}, \struck{l'action \ill{}} \add{$\Theta$} \struck{\ill{} $\Theta_{I_0}$} \add{\uncertain{laisse}} $\Pi_1(U_{0,3}, I_0)$ soit « compatible » \struck{\ill{} d'un \ill{}} \add{\ill{}} à une action $\Theta_C$ de $\Pi_1(C, I)$ \struck{qui soit l'identité \ill{}} \add{qui soit l'identité} sur les sommets $I$ (on note $\Theta_C$ \uncertain{sur} \struck{l'identité} \ill{} alors unique). On trouve ainsi une \uncertain{infinité} de conditions, \ill{} si $I \to I_0$ n'est pas surjectif) ; et dans \struck{\ill{}} \ill{} pour $I_0 \subset U_{0,3}(K)$ (NB $C_0 = U_{0,3} = \mathbb{P}^1_{\mathbb{C}} \setminus \{0, 1, \infty\}$). Pour les revêtements étales de $U_{0,3}$, on exige \add{($n \geq 1$)} quand \ill{} \ill{} il existe un $\Theta^{n}$ \add{de $\Theta$} satisfaisant aux conditions précédentes, pour $I = f^{-1}(I_0)$.

\marginal{en biais dans la marge gauche, au bas de la page : Regarder des conditions qui \ill{} \ill{} \ill{} $I_0$ \ldots}

\page{90}
\note{p. 13 de l'auteur}
Non seulement les conditions de \struck{\ill{}} « prolongeabilité » --- il \uncertain{vaut} mieux dire « d'irréductibilité » --- sous la forme précisée, \ill{} détaillée (où $\Theta$ remplacé par un $\Theta^{n}$), on trouve une condition \add{de plus}, \ill{} \ill{} un \uncertain{morphisme}
\[
  \hat{\Pi}_1(C, I) \longrightarrow \hat{\Pi}_1(X, I),
\]
que celle-ci passe au quotient --- ce qui s'exprime très exactement par la condition que pour un \ill{} $x_0 \in I$, l'action de $\Gamma$ sur
\[
  \hat{\pi}_1(C, x_0) \subset \hat{\pi}_1\bigl(U_{0,3}, f(x_0)\bigr)
\]
laisse stable le noyau de
\[
  \hat{\pi}_1(C, x_0) \longrightarrow \hat{\pi}_1(X, x_0) ;
\]
alors ce sera vrai également pour tous les autres $x_i \in I$.

Il faut exprimer la possibilité de relever l'action \struck{de} $\Theta$ à $\Pi_1(C, I)$. Notons à ce propos

\textit{Lemme}. Soit $\mathcal{C}' \to \mathcal{C}$ un hom.\ de groupoïdes connexes, qui soit injectif sur les $\pi_1$, (de sorte que $\mathcal{C}'$ est équivalent à un sous-groupoïde de $\mathcal{C}$, mais on ne suppose pas que $\mathcal{C}'$ soit lui-même un sous-groupoïde de $\mathcal{C}$, mais quand-même pas
\note{la phrase se poursuit page suivante}

\page{91}
\note{p. 14 de l'auteur ; suite du lemme commencé page précédente}
l'homomorphisme commun que
\[
  (45) \qquad \mathcal{C}' \longrightarrow \tilde{\mathcal{C}}'
\]
\note{à droite de la formule, en petit : « vrai revêtement de $\mathcal{C}$, dont les sommets sont les $\pi_0$ des fibres de $\mathcal{C}' \to \mathcal{C}$ »}
soit factorisé \ill{} de sorte qu'on peut identifier $\mathcal{C}'$ : \ill{} $\mathcal{C}' \to \mathcal{C}$ soit injectif, sous-groupoïde plein de $\tilde{\mathcal{C}}^{*}$, équivalent : $\tilde{\mathcal{C}}^{*}$ par l'inclusion, où $\tilde{\mathcal{C}}$ est un revêtement de $\mathcal{C}$.

Soit $x_0'$ un sommet de $\mathcal{C}'$, \add{$x_0$} son image dans $\tilde{\mathcal{C}}$. Soit $\Theta$ un automorphisme de $\mathcal{C}$ qui fixe $x_0$, \struck{alors} s'il \struck{\ill{}} \add{\uncertain{respecte} \ill{}} \ill{} \ill{} donnée un automorphisme \add{\uncertain{sommet}} de $\mathcal{C}$ \ill{} $\pi_1(\mathcal{C}, x_0) \overset{\mathrm{def}}{=} \Pi_0$ \ill{}, et pour tout $x_i \in \mathrm{Ob}\,\mathcal{C}$, $x_i \neq x_0$, un automorphisme des $\Pi_0(x_0)$-torseurs à droite $H_{x_i, x_0} = \mathrm{Hom}_{\mathcal{C}}(x_0, x_i)$, compatible avec l'opération de $\Theta$ sur $\Pi_0$. Ceci posé, pour qu'on \uncertain{remonte} en une action de $\tilde{\mathcal{C}}$ qui fixe $x_0'$, il faut et il suffit que le sous-groupe $\Pi_0'$ de $\Pi_0$, image de
\[
  \pi_1(\tilde{\mathcal{C}}, x_0') \hookrightarrow \pi_1(\mathcal{C}, x_0) = \Pi_0,
\]
soit \textit{stable} par $\Theta$ --- l'action de $\Theta$ sur les fibres $\tilde{\mathcal{C}}_{x_i}$ de $\tilde{\mathcal{C}}/\mathcal{C}$ est alors donnée par la formule
\[
  (45) \qquad \mathrm{Ob}\,\tilde{\mathcal{C}}_{x_i} = H_{x_i, x_0} \wedge^{\Pi_0} \Pi_0/\Pi_0'
\]
\note{sous $\Pi_0/\Pi_0'$, en petit : « ensemble des classes à droite, sur lequel $\Pi_0$ opère à g. » ; l'auteur numérote (45) deux fois sur cette page}
par transport de structure : \uncertain{partie} de \ill{}

\page{92}
\note{p. 14 de l'auteur (le numéro 14 est répété en tête de cette page)}
action sur $H_{x_i, x_0}$ le système $(\Pi_0, H_{x_i, x_0}, \Pi_0/\Pi_0')$, où on fait opérer $\Theta$ sur $\Pi_0$, $H_{x_i, x_0}$ par \struck{\ill{}} \add{la donnée}, et sur $\Pi_0/\Pi_0'$ via l'action sur $\Pi_0$, qui passe au quotient d'après la \struck{\ill{}} \add{condition} \add{\uncertain{de stabilité de $\Pi_0'$} en $\Theta$}. L'action prolongeant \struck{est} \struck{dite} ainsi \uncertain{définie} \ill{} sur $\tilde{\mathcal{C}}$ est donc unique. Pour qu'il existe une action de $\Theta^{*}$ sur $\mathcal{C}'$ qui soit compatible \ill{} $\mathcal{C}' \to \mathcal{C}$, il faut et il suffit que la condition précédente soit satisfaite, et que de \textit{plus} $\mathrm{Ob}\,\mathcal{C}' \subset \mathrm{Ob}\,\mathcal{C}$ soit stable sous \struck{\ill{}} l'action de $\tilde{\Theta}$.

Dans le cas où on ne suppose pas que l'on \struck{dispose d'}\add{donne} un $x_0' \in \mathrm{Ob}\,\mathcal{C}'$ dont on exige qu'il soit fixe par $\Theta$, on considère le \ill{} $E_0 = \tilde{\mathcal{C}}_{x_0}$, qui est un $\Pi_0$-ensemble \struck{\ill{}} homogène, et on note que les façons de \ill{} $\Theta$ en une action \struck{\ill{}} \add{$\tilde{\Theta}$} sur $\tilde{\mathcal{C}}$ correspondant \ill{} \ill{} façons de \ill{} prolonger \struck{\ill{}} ex.\ \ill{} en une action de $\tilde{\Theta}_0$ sur $E_0$, \ill{} soit compatible avec l'action sur $\Pi_0$, i.e.\ satisfasse
\[
  (46) \qquad \tilde{\Theta}_0(g \cdot x') = \Theta(g)\, \tilde{\Theta}_0(x')
\]
--- \struck{\ill{} correspond} une telle action est connue, \ill{} \ill{} \ill{} \ill{} connaît, pour \textit{un} $x_0'$ fixé de $E_0 = \mathrm{Ob}\,\tilde{\mathcal{C}}_{x_0}$, son transformé par $\tilde{\Theta}_0$, la formule (46) \struck{\ill{}} \add{$\tilde{\Theta}_0(x')$, $x' \in E_0$} donnant tous les \ill{} \ill{} \ill{}

\marginal{en biais dans la marge gauche, en face de (46) : Ceci revient à \uncertain{dire} que $\tilde{\mathcal{C}}$ est un revêtement \ill{} mais les sous-groupes \ill{} de $\Pi_0$ \ill{} \ill{} \ill{} $\Pi_0$-conjugués, \ill{} $\Pi_0$, \ill{} \ill{} \ill{} \ill{} \ill{}}

\page{93}
de $\Pi_0$, $\Pi_0' = \mathrm{Im}\bigl(\pi_1(\tilde{\mathcal{C}}, x_0') \to \pi_1(\mathcal{C}, x_0)\bigr)$, de sorte que
\[
  (47) \qquad E_0 \simeq \Pi_0/\Pi_0' ;
\]
la condition d'existence de $\tilde{\Theta}_0$ revient à dire que $\Theta(\Pi_0')$ est conjugué à $\Pi_0'$ dans $\Pi_0$, et les choix possibles pour \uncertain{remonter} $\Theta$ en $\tilde{\Theta}_0$, donc pour remonter $\Theta$ en $\tilde{\Theta}$, correspondent aux éléments $x''$ de $E_0$ \struck{tels que} \add{de la forme} $g(x_0')$ où $g \in \Pi_0$ satisfait
\[
  (48) \qquad \mathrm{int}\, g\,(\Pi_0') = \Theta(\Pi_0'),
\]
i.e.\ \uncertain{tels} que $g \in \mathrm{Transp}_{\Pi_0}\bigl(\Pi_0', \Theta(\Pi_0')\bigr)$. \add{Cette condition ne dépend que de l'action \textit{sur} $\Pi_0'$ définie par $\Theta$ sur $\Pi_0'$ --- mais les façons de remonter $\Theta$ en $\tilde{\Theta}$, \uncertain{elles}, en dépendent, au \ill{} de \ill{} \ill{} près. \ill{} $\Theta$ \ill{}}

Donc les $\tilde{\Theta}$ possibles correspondant aux éléments de
\[
  (49) \qquad \mathrm{Transp}_{\Pi_0}\bigl(\Pi_0', \Theta(\Pi_0')\bigr)/\Pi_0'
\]
\note{sous (49), en petit : « ensemble homogène à droite ($\Pi_0'$ opère à droite sur le transporteur, via $g \mapsto g\alpha$ pour $\alpha \in \Pi_0'$) »}

Ainsi, on voit que dans le cas
\[
  \mathcal{C} = \hat{\Pi}_1(X_{\bar K}, I), \qquad I \subset X_K,
\]
\struck{\ill{}} quand $\mathcal{C}'$ \struck{\ill{} provient} est de la forme $\hat{\Pi}_1(X', I')$, où $X'$ est un revêtement étale (fini) de $X$, \add{muni d'un} $I' \to I$, \add{la condition qu'un} $\Theta$ se remonte \add{à $\mathcal{C}'$} \ill{} que la condition \ill{} sur action de $\mathcal{C}'$ est une condition algébrique des plus
\note{la phrase se poursuit page suivante}

\marginal{en biais dans la marge gauche, en haut : En effet, il faut que $\mathrm{Im}\bigl(\pi_1(\tilde{\mathcal{C}}, x_0') \to \pi_1(\mathcal{C}, x_0)\bigr)$ soit $\Theta(\Pi_0')$ \ill{} $= \pi_1(\ldots)$, $g \in \Pi_0$, \ill{} $x_1' = \tilde{\Theta}(x_0')$, si on \ill{} $x_1' = g(x_0')$, cet \ill{}}

\marginal{en biais dans la marge gauche, en bas, sous un long passage biffé illisible : Bien noter que \ill{} \ill{} \ill{} $\Theta$ \ill{} \ill{} \ill{} \ill{} $\mathcal{C}' \subset \tilde{\mathcal{C}}$ \ill{} $\Theta$ \ill{} \ill{}}
\note{un passage marginal plus long, entre ces deux notes, est entièrement biffé de traits croisés et illisible ; on y distingue « (48) » et « $\Theta(\Pi_0')$ »}

\page{94}
\note{p. 15 de l'auteur}
\struck{plus} simple, que pour $I' = f^{-1}(I)$ ne dépend pas du choix de $I'$ mais de $I$. Si on veut que l'action remontée $\Theta'$ soit l'identité sur $I'$, on a unicité, et la condition supplémentaire \uncertain{signifie} bien que $\tilde{\Theta}$ doit bien opérer trivialement sur $I'$.

Résumons : l'action de $\Gamma$ \add{$= \mathrm{Gal}\,\bar{K}/K$} sur un $\mathcal{C} = \hat{\Pi}_1(X, I)$ avec $I \subset X_K$, on trouve des \add{conditions nécessaires qu'une} condition \ill{} \add{\uncertain{satisfait} sur \ill{} \ill{}} action $\Theta$ de $\mathcal{C}$ \add{provienne d'un} $u \in \Gamma$ :
\begin{enumerate}
  \item[a)] Pour tout revêtement fini \add{étale} $X'$ de $X$ \add{(défini sur $K)$}, \struck{\ill{} considérons} \struck{il y a l'action de} \add{une action} $\tilde{\Theta}$ \add{de $\Theta$} sur $\hat{\Pi}_1(X', I')$, où $I' = X' \mid I$, qui remonte $\Theta$, et qui \add{soit} l'identité sur $I' \cap X'_K$.
  \item[b)] De plus, pour \ill{} \add{\uncertain{tout} \ill{} (défini sur $K$)} $X' \xrightarrow{\;f\;} Y$, et tout $x' \in I' \cap X'_K$, le noyau de
  \[
    (*) \qquad \hat{\pi}_1(X', x') \longrightarrow \hat{\pi}_1\bigl(Y, f(x')\bigr)
  \]
  \add{(pour $\tilde{\Theta}$ \uncertain{convenable} !)} est stable sous $\tilde{\Theta}^{*}$ \add{--- plus généralement,} sans supposer $x' \in X'_K$, mais choisissant $\nu \in \mathbb{N}^{*}$ tel que $\Theta^{\nu} x' = x'$, il faut que \struck{\ill{}} le noyau de $(*)$ soit stable par $\Theta^{\nu}$.
\end{enumerate}
\note{dans (*), le premier argument de $\hat{\pi}_1$ au but est surchargé ; on lit $Y$}

Ce sont là des conditions nécessaires. Je pense qu'on \ill{} pas trop \ill{},

\marginal{en biais dans la marge gauche, séparé du texte par un trait vertical : Cas d'une action $\Theta$ \ill{} sur $\mathcal{C}$ qui \uncertain{fixe} \ill{} $x_0 \in I$ \ill{} \ill{} \ill{} \ill{} condition \ill{} \ill{} \ill{} \ill{}. NB Ces conditions, \ill{}, \ill{} \ill{} \ill{} \ill{} \ill{} \ill{} \ill{} \ill{} \uncertain{définie} \ill{} de $\Gamma$ \ill{} \ill{}}

\page{95}
\uncertain{exagère} en conjecturant que \uncertain{en général}, ces conditions, \add{nécessaires}, sont aussi \ill{} suffisantes \struck{pour} (même \ill{} la considération des $\Theta^{\nu}$) \ill{} \ill{} pour que $\Theta$ provienne bien \ill{} d'un $u \in \Gamma$, \struck{\ill{}} lorsque $X$ est (disons) une courbe algébrique anabélienne, avec $K$ extension de type fini de $\mathbb{Q}$. Bien \ill{} la question de l'unicité de $u$ ne se pose pas ici --- quand $K$ est \uncertain{une extension} algébrique de $\mathbb{Q}$, on s'attend à une unicité de $u$. \uncertain{Quand} $K$ est de degré de transcendance $> 0$, dès que \ill{}, je pense pouvoir s'attendre : l'unicité de $u$, il faut au moins supposer que le corps \uncertain{modulaire} de $X$ a même degré de transcendance que $K$, i.e.\ que $K$ est une extension finie de ce corps modulaire. Je ne suis pas étonné qu'on ait unicité dans ce cas.
\note{« Bien » ouvre, après « $\mathbb{Q}$. », un passage que l'auteur entoure d'un trait anguleux}

Je reviens sur la question de la relevabilité d'un automorphisme $\Theta$ d'un groupoïde connexe $\mathcal{C}$ : un revêtement $\tilde{\mathcal{C}}$, quand $\Theta$ fixe un sommet

\marginal{en biais dans la marge gauche, en haut : Je ne vois pas très bien que les conditions a), b), \ill{} \ill{} \ill{} \ill{} \ill{} dans \ill{} \ill{} \ill{} \ill{} de pro-\ill{} \ill{} $\pi$ \ill{} $\Theta$ \ill{} \ill{} \ill{} \ill{} \ill{} \ill{} \ill{} \ill{} \ill{} CN \ill{} \ill{} \ill{} \ill{} \ill{} \ill{} \ill{} $u \in \Gamma$ \ill{}}

\marginal{en biais dans la marge gauche, plus bas, en partie encadré : Le \ill{} \ill{} \ill{} \ill{} \ill{} \ill{} \ill{} \ill{} \ill{} \ill{} \ill{} \ill{} \ill{} \ill{} \ill{} \ill{} Belyi \ill{} \ill{} \ill{} \ill{} \ill{} \ill{} \ill{} \ill{} condition \ill{} \ill{} \ill{} \ill{}}
\note{les deux notes marginales de cette page, écrites en biais et serrées, sont en grande partie illisibles}

\page{96}
\note{p. 16 de l'auteur}
$x_0$ de $\mathcal{C}$, mais sans supposer maintenant $\tilde{\mathcal{C}}$ connexe. (\struck{\ill{}} Revêtement \ill{} à $\tilde{\mathcal{C}}$ proviendrait d'une situation géométrique $X' \to X$, avec $X$ connexe, \struck{et \ill{}} mais $X'$ \struck{\ill{}} seulement $K$-connexe i.e.\ connexe en tant que schéma sur $K$ --- ce qui signifie qu'il peut se décrire en termes de $E_0 = X'_{x_0}$, sur lequel opèrent \uncertain{directement} $\Gamma \cdot \Pi_0$ ($\Pi_0 = \pi_1(X, x_0) = \pi_1(\mathcal{C}, x_0)$) \uncertain{transitivement}, mais sans que l'opération de $\Pi_0$ soit transitive.) Associons à tout $x' \in E_0$ son stabilisateur dans $\Pi_0$, on trouve une application
\[
  (50) \qquad
  \begin{aligned}
    E_0 &\longrightarrow \text{Ens.\ des s-groupes de } \Pi, \\
    x' &\longmapsto \text{stabilisateur } \Pi_{x'},
  \end{aligned}
\]
qui commute aux actions de $\Pi_0$ (cop.\ sur le 2\textsuperscript{e} membre par automorphismes intérieurs), d'où par passage aux quotients
\[
  (51) \qquad E_0/\Pi_0 \overset{\mathrm{def}}{=} \pi_0(E_0) \longrightarrow \Sigma = \mathrm{Ssgr}(\Pi)/\Pi
\]
\note{sous $\pi_0(E_0)$, en petit : « i.e.\ les $\pi_0$ des \uncertain{fibres} des $\Pi$-ens. » ; à droite, un symbole surchargé, lu $\mathrm{Ssgr}$ (ensemble des sous-groupes)}

D'autre part, $\Theta$ opère sur le 2\textsuperscript{ième} membre de (51) par transport de structure, via $\Pi$, et \ill{} une condition nécessaire

\page{97}
évidente pour que $\Theta$ se relève en un $\tilde{\Theta}$, c'est que l'image de l'application (51) soit \textit{stable} sous l'action de $\Theta$ ; en fait, l'action $\tilde{\Theta}$ présumée, \struck{\ill{}} \add{\uncertain{montrerait} à $E_0$, \uncertain{doit donner}} \add{On voit donc qu'un $\tilde{\Theta}$ \ill{} deux pts de $\Sigma$ \ill{}} \struck{\ill{} \ill{}} la compatibilité \add{\uncertain{de cette action} $\Theta$,} se déduit de (52) \struck{\ill{}} \add{\uncertain{(considéré} \ill{} \ill{} \ill{}) \uncertain{dont}} \struck{\ill{} stabilité} \struck{\ill{}} que $\Gamma$ vienne d'\ill{} signifie justement qu'il existe une \struck{\ill{}} \add{bijection}
\[
  (52) \qquad \pi_0(\tilde{\Theta}_0) : \pi_0(E_0) \xrightarrow{\;\sim\;} \pi_0(E_0)
\]
\struck{stable que} qui soit compatible avec l'action de $\Theta$ sur le deuxième membre de (51). Ceci dit, pour toute telle relèvement, il existe \struck{\ill{}} \add{\ill{}} \textit{tel} \textit{de} \uncertain{relever} $\tilde{\Theta}_0 : E_0 \to E_0$ \add{compatible avec les actions de $\Pi$ et l'action de $\Theta$ sur $\Pi$ (cf.\ (46))}, \uncertain{si} on connaît \uncertain{quand} on lorsque $\tilde{\Theta}$ est \ill{} \ill{} connaît son \ill{} \uncertain{valeurs} $\tilde{\Theta}(x_i)$ \struck{\ill{}} pour \textit{un} $x_i \in$ \struck{\ill{}} \add{$E_0$} dans chaque des \add{\uncertain{composantes}} composantes connexes de $E_0$ \ldots
\note{la page est très corrigée ; les ajouts interlinéaires s'enchevêtrent avec les biffures et l'ordre de lecture proposé ici est incertain}

Pour les automorphismes $\Theta$ \uncertain{de l'un} $\hat{\Pi}_1(X, I)$, $I \subset X_K$, pour \uncertain{provenir} de $\Gamma$, j'ai \ill{} \uncertain{cependant} \ill{} plus \uncertain{importantes} des conditions nécessaires, \ill{} \ill{} \ill{} \ill{} : \ill{} \ill{} \uncertain{fibres} \ill{} \ill{} \ill{} !
\note{la fin de cette page, d'une écriture très rapide, n'est lue qu'en partie}

\page{98}
\note{p. 17 de l'auteur}
C'est la condition de compatibilité avec la « structure : \uncertain{bout} » --- dans le cas où $X$ est une courbe algébrique, disons (\uncertain{sinon}, on sent bien qu'il faut introduire, \uncertain{dans} les cas généraux, une « structure : l'infini » sur un groupoïde fondamental \ldots). Cette compatibilité \add{justement} s'exprime de façon particulièrement \ill{} et élégante dans le contexte des groupoïdes, \ill{} \ill{} \ill{} compatibilité \ill{} la \uncertain{notion} d'une \ill{} de \ill{} (cf.\ (39)) \ill{} \uncertain{solutions}, \ill{} \ill{} \ill{} \ill{} \uncertain{solutions} \ill{} \ill{} \ill{}, \ill{} \ill{} \ill{} « \ill{} » \ill{} --- \ill{} \ill{} \ill{} sur ce point de vue \uncertain{plus} bas.

\struck{Remarques.} Revenant aux conditions a) b), il est clair \uncertain{d'emblée} qu'elles sont \uncertain{inexplicitables} (c.à.d.\ inexprimables) en termes $\pm$ directs, comme des conditions \uncertain{portant directement} sur \uncertain{des} « paramètres » \struck{\ill{}} \add{tels que} du type $g_{\ell}(u)$ (des paramètres $\alpha, \beta, \xi, \eta$ dont nous \ill{} \ill{} \ill{} des automorphismes : l'action de $\Pi_{0,3}$ \ldots) qui « repèrent » un $u \in \Gamma$ (ici $\Gamma = \mathrm{Gal}\,\bar{\mathbb{Q}}/\mathbb{Q}$). \struck{Probablement} C'est pour cette raison que le
\note{la phrase se poursuit page suivante ; page d'une écriture rapide, nombreuses lectures douteuses}

\page{99}
\ill{} plus haut (cf.\ \uncertain{une remarque}) \uncertain{sur} les conditions des lacets, \uncertain{soit} \ill{} n'est pas bien intéressante, en \ill{} pratique. Mais je \ill{} que \uncertain{vu} l'étude des $\mathbb{M} \Delta$ \ill{} $M_{g,\nu}$ et des groupes de Teichmüller, je suis en train de \ill{} \ill{} \ill{} une conjecture \ill{} \ill{} \ill{} intéressante et plausible, en termes de propriétés \ill{} des groupes profinis \uncertain{tels} que $\hat{\Pi}_{0,3}$, \ill{} par \uncertain{cohomologie} \ill{} \uncertain{(on} \ill{} \ill{} par des conditions liées \ill{} $M_{0,4}$, $M_{0,5}$ \ldots).

Je pense en venir encore une fois, sous forme complétée, les conditions pour un automorphisme $\Theta$ \add{d'un \uncertain{groupoïde}} $\hat{\Pi}_1(X, I)$, $I \subset X_{\bar K}$, \add{$I \cap X_K \neq \emptyset$} (je pense \uncertain{surtout}, \ill{} au cas $X = U_{0,3}$) \uncertain{provienne} d'un $u \in \Gamma = \mathrm{Gal}\,\bar{K}/K$ :

1) Conditions des lacets (pour mémoire --- c'est sûrement la plus importante de toutes, \uncertain{et de} très loin !)

2) Condition de stabilité pour revêtements étales $X'$ de $X$ (définis sur $K$, bien sûr) :
\note{la phrase se poursuit au-delà de cette page}

\marginal{en biais dans la marge gauche, en haut : En \ill{} \ill{} \ill{} \ill{} conjecture \ill{} \ill{} \ill{} \ill{} \ill{} $\mathrm{Gal}\,\bar{\mathbb{Q}}/\mathbb{Q}$, $I$ \ill{} \ill{} \ill{} \ill{} \ill{} \ill{} \ill{} \ill{} \ill{} \ill{} $X$ \ill{} \ill{}}
\note{la note marginale, écrite en biais et en partie barrée d'un trait, est presque entièrement illisible}

\page{100}
\note{p. 18 de l'auteur ; suite de la condition 2) de la page précédente}
\ill{} de $\Theta$ se relève en une action $\Theta'$ de $\Theta$ sur $\hat{\Pi}_1(X', I')$, où $I' = X' \mid I$. \struck{\ill{}} Si $I' \cap X'_K \neq \emptyset$, $X'$ étant connexe (il suffit que $\Gamma$ soit transitif sur $\pi_0(X')$ i.e.\ que $X'$ soit « $K$-connexe » i.e.\ connexe en tant que schéma sur $K$) donc l'action \uncertain{relevée} est unique.

\add{La condition 2 s'explicite en pratique en une condition de stabilité de \uncertain{certains} sous-groupes ouverts d'un $\pi_1(X, x)$ ($x \in X_K \cap I$) \ill{} \ill{} \ill{} $\Theta$ \ill{}}

3) Pour les \ill{} homom.\ $f : X' \to Y$ et \struck{\ill{}} \add{\uncertain{tout}} $x' \in X'_K \cap I'$, \struck{l'action} la $\Theta'$ \struck{\ill{}} de \uncertain{telle} \uncertain{action} \uncertain{doit} \ill{} \ill{} de $\hat{\pi}_1(X', x')$ qu'il \ill{} \uncertain{stabilise} \uncertain{laisse} stable le noyau de
\[
  \hat{\pi}_1(X', x) \longrightarrow \hat{\pi}_1\bigl(Y, f(x)\bigr) .
\]
C'est encore une condition de stabilité pour certains sous-groupes, mais \uncertain{pour une deuxième fois} \ill{} $\pi_1(X, x)$ \ill{} $\Theta$.

4) Soient $X'_1$, $X'_2$ deux rev.\ étales finis de $X$, et $f$ un morphisme $X'_1 \xrightarrow{\;f\;} X'_2$ (\uncertain{pas} \uncertain{nécessairement} un $X$-morphisme !), $I'_1 = X'_1 \mid I$, $I'_2 = X'_2 \mid I$ ; on suppose \struck{plus précisément que} $I'_1 \cap f^{-1}(I'_2) \neq \emptyset$ ; pour $J \overset{\mathrm{def}}{=} I'_1 \cap f^{-1}(I'_2)$, alors \struck{\ill{}} le morphisme
\[
  \text{\struck{$I'_1 \cap f^{-1}(I'_2) \longrightarrow$}} \quad
  \hat{\Pi}_1(X'_1, J) \longrightarrow \hat{\Pi}_1(X'_2, I'_2)
\]
est compatible \add{\uncertain{aux}} \add{\uncertain{prolongements}} \ill{} actions de $\Theta$ sur ces groupoïdes, s'il y \ill{} (pour \ill{} \ill{} \ill{} \ill{} \uncertain{compatibilité}, cf.\ 2.).
\note{l'argument se poursuit au-delà de ce lot}

\marginal{en biais dans la marge gauche, en haut : \ill{} \ill{} \ill{} \ill{} \ill{} $I \cap X'_K$ \ill{} \ill{} \ill{} $\Theta'$ \ill{} $\Theta$ \ill{} \ill{} $X'_K \neq \emptyset$ \ill{}}

\marginal{en biais dans la marge gauche, plus bas : \ill{} \ill{} \ill{} \ill{} $x' \in X'_K \cap I'$, \ill{} \ill{} \ill{} $\Theta_{\nu}$ \ill{} \ill{} \ill{} $\hat{\pi}_1(X', x')$ \ill{} $\Theta_{\nu} x' = x'$ \ill{}}
\note{les deux notes marginales de cette page, serrées et en biais, sont presque entièrement illisibles}
\end{document}
